\documentclass[a4paper,14pt]{extarticle}

\usepackage{fontspec}
\usepackage{polyglossia}
\setmainlanguage{russian}
\setotherlanguage{english}
\setmainfont{Times New Roman}
\setsansfont{Arial}
\setmonofont{Consolas}[Script=Cyrillic]
\newfontfamily\cyrillicfonttt{Consolas}[Script=Cyrillic]

\usepackage[left=20mm,right=15mm,top=15mm,bottom=15mm]{geometry}
\usepackage{setspace}
\usepackage{indentfirst}
\usepackage{titlesec}
\usepackage{tocloft}
\usepackage{caption}
\usepackage{array}
\usepackage{longtable}
\usepackage{booktabs}
\usepackage{tabularx}
\usepackage{enumitem}
\usepackage{float}
\usepackage{placeins}
\usepackage{lastpage}
\usepackage{listings}
\usepackage{xcolor}
\usepackage{xurl}
\usepackage{hyperref}

\hypersetup{
    colorlinks=false,
    pdfborder={0 0 0},
    pdftitle={Пояснительная записка QuasarDB},
    pdfauthor={Команда QuasarDB}
}

\onehalfspacing
\setstretch{1.15}
\setlength{\parindent}{1.25cm}
\setlength{\parskip}{0pt}
\emergencystretch=3em
\sloppy

\titleformat{\section}
  {\normalfont\bfseries\fontsize{16pt}{24pt}\selectfont}
  {\thesection}{1em}{}
\titleformat{\subsection}
  {\normalfont\bfseries\fontsize{16pt}{24pt}\selectfont}
  {\thesubsection}{1em}{}
\titlespacing*{\section}{0pt}{1.5\baselineskip}{0.8\baselineskip}
\titlespacing*{\subsection}{0pt}{1.2\baselineskip}{0.6\baselineskip}

\renewcommand{\cftsecleader}{\cftdotfill{\cftdotsep}}
\renewcommand{\contentsname}{Содержание}
\addto\captionsrussian{%
  \renewcommand{\figurename}{Рисунок}%
  \renewcommand{\tablename}{Таблица}%
  \renewcommand{\lstlistingname}{Листинг}%
}

\DeclareCaptionFont{captiontwelve}{\fontsize{12pt}{12pt}\selectfont}
\captionsetup{
    font={captiontwelve,it},
    labelfont={captiontwelve,it},
    labelsep=period,
    justification=centering,
    singlelinecheck=false
}
\captionsetup[figure]{position=bottom,justification=centering}
\captionsetup[table]{position=top,justification=raggedright,singlelinecheck=false}
\captionsetup[lstlisting]{position=top,justification=raggedright,singlelinecheck=false,labelsep=period}
\urlstyle{same}
\Urlmuskip=0mu plus 2mu

\lstdefinestyle{qdbcpp}{
    basicstyle=\ttfamily\fontsize{12pt}{12pt}\selectfont,
    keywordstyle=\color{blue!60!black},
    commentstyle=\color{gray!70!black},
    stringstyle=\color{green!40!black},
    numbers=left,
    numberstyle=\tiny,
    stepnumber=1,
    numbersep=8pt,
    showstringspaces=false,
    breaklines=true,
    frame=single,
    tabsize=4,
    captionpos=t
}

\pagestyle{plain}
\pagenumbering{arabic}

\newcommand{\University}{Московский авиационный институт}
\newcommand{\Department}{Кафедра системного программирования}
\newcommand{\WorkTitle}{Разработка клиент-серверной системы управления базами данных QuasarDB}
\newcommand{\CourseName}{Курсовая работа по дисциплине «Системное программирование»}
\newcommand{\Team}{Эльдар С., Кирилл П., Михаил С., Кирилл Б.}
\newcommand{\Group}{учебная группа: \underline{\hspace{4cm}}}
\newcommand{\Supervisor}{руководитель: \underline{\hspace{4cm}}}
\newcommand{\breaktt}[1]{\texttt{\footnotesize #1}}
\newcolumntype{Y}{>{\raggedright\arraybackslash}X}

\begin{document}

\begin{titlepage}
    \thispagestyle{empty}
    \begin{center}
        Министерство науки и высшего образования Российской Федерации\\
        Федеральное государственное автономное образовательное учреждение\\
        высшего образования\\
        «\University»\\
        \Department

        \vfill

        \textbf{\CourseName}\\[1.5cm]
        \textbf{\WorkTitle}\\[0.5cm]
        Пояснительная записка

        \vfill

        \begin{flushright}
            Выполнили: \Team\\
            \Group\\[0.5cm]
            Проверил: \Supervisor
        \end{flushright}

        \vfill

        Москва, 2026
    \end{center}
\end{titlepage}

\clearpage
\section*{Реферат}
\addcontentsline{toc}{section}{Реферат}

Пояснительная записка посвящена разработке QuasarDB — учебной клиент-серверной системы управления базами данных на языке C++20. В работе описаны требования к системе, архитектура комплекса приложений, SQL-подобный язык запросов, организация хранения данных на диске, индексирование на основе B*+-дерева, журналирование изменений, механизм темпорального отката, сетевой протокол, аутентификация, разграничение доступа, логирование и телеметрия.

Цель работы — реализовать СУБД, поддерживающую целочисленные и строковые типы данных, двухуровневую иерархию баз данных и таблиц, интерактивный и пакетный режимы работы, синтаксическую и семантическую валидацию запросов, JSON-формат результата и оптимизацию выборок с использованием индексов.

\clearpage
\tableofcontents
\clearpage

\section*{Перечень сокращений и обозначений}
\addcontentsline{toc}{section}{Перечень сокращений и обозначений}

\begin{longtable}{>{\raggedright\arraybackslash}p{0.22\textwidth}>{\raggedright\arraybackslash}p{0.70\textwidth}}
\caption{Перечень сокращений и обозначений}\\
\toprule
Сокращение & Расшифровка \\
\midrule
\endfirsthead
\caption*{Продолжение таблицы \thetable}\\
\toprule
Сокращение & Расшифровка \\
\midrule
\endhead
AST & абстрактное синтаксическое дерево запроса. \\
DDL & подмножество языка для определения структуры данных. \\
DML & подмножество языка для манипулирования данными. \\
IPC & межпроцессное взаимодействие. \\
JSON & текстовый формат обмена структурированными данными. \\
JWT & подписанный токен для аутентификации и авторизации. \\
RBAC & модель разграничения доступа на основе ролей и прав. \\
WAL & журнал предварительной записи. \\
СУБД & система управления базами данных. \\
\bottomrule
\end{longtable}

\clearpage
\section{Введение}

В рамках курсовой работы разработана система управления базами данных QuasarDB. Проект выполняется как программный комплекс из нескольких приложений: клиентского терминала, входного серверного узла и узлов хранения. Такая декомпозиция позволяет отделить пользовательский интерфейс и сетевой протокол от физического хранения данных, а также подготовить систему к распределенной обработке запросов.

Актуальность работы связана с необходимостью практического изучения системного программирования: работы с файловой системой, сетевыми сокетами, потоками, сериализацией данных, структурами индексации и механизмами восстановления состояния. В отличие от простого интерпретатора запросов, QuasarDB хранит данные на диске, ведет журналы изменений, поддерживает индексы и возвращает результаты выборок в формате JSON.

Основными требованиями к проекту являются: реализация на языке C++ стандарта C++14 и выше, поддержка типов \texttt{int} и \texttt{string}, хранение данных в файловой системе без потери после перезапуска, SQL-подобный язык запросов, синтаксическая и семантическая валидация, интерактивный и пакетный режимы клиента, индексирование полей с модификатором \texttt{INDEXED}, запрет дублирования данных в индексах и информативная обработка ошибок.

Цель пояснительной записки — описать принятые проектные решения, используемые средства языка C++, алгоритмы и структуры данных для реализованных заданий курсового проекта.

\section{Постановка задачи}

Требуется реализовать СУБД, в которой уровень системы содержит набор баз данных, база данных содержит набор таблиц, а таблица хранит типизированные записи. Взаимодействие с системой осуществляется через SQL-подобный язык. Команды завершаются символом \texttt{;} и могут занимать несколько строк. Имена сущностей состоят из латинских букв, цифр и символа подчеркивания, но не начинаются с цифры.

Система должна поддерживать операции создания и удаления баз данных, выбора активной базы, создания и удаления таблиц, вставки, обновления, удаления и выборки записей. Для условий выборки используются сравнения \texttt{==}, \texttt{!=}, \texttt{<}, \texttt{>}, \texttt{<=}, \texttt{>=}, диапазонный предикат \texttt{BETWEEN}, регулярное сравнение \texttt{LIKE}, а также составные логические выражения \texttt{AND} и \texttt{OR}.

Дополнительные задания проекта включают темпоральную персистентность без полных снимков файлов, интернирование строк, клиент-серверную и кластерную архитектуру, heartbeat-проверки узлов хранения, асинхронную обработку длительных операций, access logs, телеметрию, RBAC-аутентификацию, значения по умолчанию и агрегатные функции \texttt{SUM}, \texttt{COUNT}, \texttt{AVG}.

\begin{table}[H]
\caption{Соответствие требований и подсистем QuasarDB}
\centering
\small
\begin{tabularx}{\textwidth}{>{\raggedright\arraybackslash}p{0.26\textwidth}Y}
\toprule
Требование & Реализация в проекте \\
\midrule
Хранение данных & Файлы \texttt{.data}, \texttt{.schema}, \texttt{.bin}, \texttt{.jnl}, \texttt{.idx} в каталоге таблицы. \\
Индексы & B*+-дерево хранит ключи и адреса записей, сами записи не дублируются. \\
SQL-подобный язык & Лексер, парсер, AST, грамматика в документации проекта. \\
Клиент-сервер & CLI-клиент, Entrypoint-сервер, Storage-узлы, общий сетевой слой. \\
Надежность & Страничная запись, синхронный журнал обратных операций, восстановление через \texttt{REVERT}. \\
Безопасность & Хранилище аккаунтов, PBKDF2-HMAC-SHA256, JWT, проверка RBAC. \\
\bottomrule
\end{tabularx}
\end{table}

\begin{table}[H]
\caption{Покрытие дополнительных заданий}
\centering
\small
\begin{tabularx}{\textwidth}{>{\raggedright\arraybackslash}p{0.12\textwidth}>{\raggedright\arraybackslash}p{0.31\textwidth}Y}
\toprule
Номер & Задание & Реализация \\
\midrule
1 & Темпоральная персистентность & Журнал обратных операций \texttt{Journal}; команда \texttt{REVERT} применяет записи журнала от конца к заданному времени. \\
2 & Интернирование строк & \texttt{Interner} хранит уникальные строки в \texttt{std::deque} и сопоставляет их через \texttt{std::unordered\_map}. \\
3 & Клиент-сервер & CLI-клиент передает JSON-запросы Entrypoint-серверу по TCP. \\
4 & Кластер & Entrypoint, \texttt{Registry}, \texttt{Catalog} и Storage-процессы разделяют маршрутизацию и физическое хранение. \\
5 & Heartbeat & \texttt{Monitor} опрашивает Storage через \texttt{ping} и перезапускает недоступные процессы. \\
6 & Асинхронные задачи & \texttt{TaskProcessor} возвращает GUID v4 и хранит статус до запроса \texttt{check\_task}. \\
7 & Access logs & \texttt{Logger} фиксирует пользователя, обработчик, тело запроса, длительность и статус. \\
8 & Телеметрия & \texttt{Telemetry} поддерживает RPS, окна 10 минут, 10 секунд и 60 секунд. \\
9 & RBAC и JWT & \texttt{AccountStore}, \texttt{JwtHandler} и \texttt{RBACManager} проверяют пользователя и права. \\
10 & \texttt{DEFAULT} & \texttt{Schema} и \texttt{Table::apply\_defaults} подставляют значения при \texttt{INSERT}. \\
11 & \texttt{AND}/\texttt{OR} & Парсер строит AST с приоритетом \texttt{AND} выше \texttt{OR}, исполнитель вычисляет трехзначную SQL-логику. \\
12 & Агрегаты & \texttt{Executor::AggregateSelect} реализует \texttt{SUM}, \texttt{COUNT}, \texttt{AVG}. \\
\bottomrule
\end{tabularx}
\end{table}

\FloatBarrier

\section{Архитектура программного комплекса}

QuasarDB построена как набор CMake-таргетов. Общие сетевые и вспомогательные классы вынесены в библиотеку \texttt{libs/core}; клиент расположен в \texttt{apps/client}; входной сервер — в \texttt{apps/server}; узел хранения — в \texttt{apps/storage}. Такое разделение уменьшает связность компонентов и позволяет тестировать подсистемы независимо.

\begin{table}[H]
\caption{Назначение основных модулей}
\centering
\small
\begin{tabularx}{\textwidth}{>{\raggedright\arraybackslash}p{0.24\textwidth}Y}
\toprule
Модуль & Ответственность \\
\midrule
\texttt{libs/core} & Содержит TCP-клиент, TCP-сервер, абстракции сокетов, JSON-ответы, потоковый пул и утилиты работы с JSON-файлами. \\
\texttt{apps/client} & Отвечает за чтение команд, пользовательскую сессию, отправку запросов, получение ответов и форматированный вывод результата. \\
\texttt{apps/server} & Реализует публичный API, лексер, парсер, анализатор, маршрутизатор, каталог, регистрацию Storage-узлов, безопасность, логи и метрики. \\
\texttt{apps/storage} & Выполняет операции над таблицами, страницами, схемами, индексами, журналом и пулом строк. \\
\bottomrule
\end{tabularx}
\end{table}

\begin{figure}[H]
\centering
\fbox{\begin{minipage}{0.92\textwidth}
\centering
\texttt{qdb-client} $\rightarrow$ \texttt{qdb-server / Entrypoint} $\rightarrow$ \texttt{qdb-storage}\\[0.4cm]
\texttt{libs/core}: TCP-сокеты, JSON-кадры, общие утилиты
\end{minipage}}
\caption{Общая схема взаимодействия компонентов}
\end{figure}

Клиентский терминал читает команды пользователя, накапливает многострочный ввод до символа завершения, отправляет запросы на сервер и форматирует JSON-ответы. Entrypoint принимает клиентские соединения, выполняет разбор SQL, семантические проверки, проверку прав доступа, маршрутизацию и сбор служебных метрик. Storage получает не текст SQL, а сериализованное AST-представление и выполняет операции над локальными файлами таблицы.

Жизненный цикл запроса состоит из нескольких этапов. На первом этапе клиент формирует JSON-сообщение с действием \texttt{query}. На втором этапе Entrypoint проверяет токен, разбирает SQL-строку и строит AST. На третьем этапе анализатор определяет активную базу данных, таблицу, требуемые права и типы выражений. На четвертом этапе маршрутизатор выбирает Storage-узел, а Storage выполняет команду и возвращает JSON-результат. Такой пайплайн отделяет ошибки синтаксиса, ошибки прав доступа и ошибки физического исполнения.

\subsection{Используемые средства C++}

Проект написан на C++20, но не привязан к экспериментальным расширениям компилятора. В коде активно используются стандартная библиотека, RAII, умные указатели, move-семантика, шаблоны, исключения и потоковые примитивы. Это важно для системного проекта: сетевые соединения, файловые дескрипторы, буферы страниц и фоновые потоки должны освобождаться предсказуемо даже при ошибках.

\begin{table}[H]
\caption{Средства C++ и их роль в проекте}
\centering
\small
\begin{tabularx}{\textwidth}{>{\raggedright\arraybackslash}p{0.30\textwidth}Y}
\toprule
Средство & Использование \\
\midrule
\texttt{std::unique\_ptr}, RAII & Владение серверами, узлами дерева, клиентскими соединениями и автоматическое закрытие ресурсов в деструкторах. \\
\texttt{std::variant} & Представление значений \texttt{NULL}, \texttt{INT}, \texttt{STRING} в едином типе \texttt{Value}. \\
\texttt{std::filesystem} & Создание каталогов баз данных, путей таблиц, файлов схем, журналов и индексов. \\
\texttt{std::thread}, атомарные счетчики & Обработка клиентских подключений, heartbeat-мониторинг, фоновые задачи и телеметрия. \\
Шаблоны и concepts & Обобщенная реализация B*+-дерева для целочисленных, строковых и адресных ключей. \\
\texttt{std::optional} & Возврат отсутствующих записей, токенов, задач и необязательных частей AST без специальных sentinel-значений. \\
\texttt{nlohmann::json} & Единый формат публичного и внутреннего сетевого API, результатов выборок и конфигурационных файлов. \\
\bottomrule
\end{tabularx}
\end{table}

\subsection{Сетевой протокол}

Обмен сообщениями выполняется поверх TCP. Каждое сообщение кодируется как JSON в UTF-8 и передается с префиксом длины: сначала четыре байта размера полезной нагрузки в сетевом порядке байтов, затем JSON-тело. Такой формат отделяет транспортный уровень от содержимого запроса и позволяет корректно читать сообщения независимо от фрагментации TCP-потока.

Публичный запрос содержит поля \texttt{action}, \texttt{token} и \texttt{data}. Ответ содержит \texttt{status}, \texttt{message} и \texttt{data}. Статус принимает значения \texttt{success}, \texttt{error} или \texttt{pending}. Для длительных операций сервер возвращает \texttt{task\_id}, а клиент может позднее запросить состояние через действие \texttt{check\_task}.

\begin{lstlisting}[style=qdbcpp,caption={Пример JSON-запроса к публичному API}]
{
  "action": "query",
  "token": "JWT_TOKEN",
  "data": {
    "query": "SELECT COUNT(id) AS total FROM users WHERE age >= 18;"
  }
}
\end{lstlisting}

\subsection{Кластерная модель}

Серверная часть разделена на Entrypoint и Storage. Entrypoint играет роль многопоточного балансировщика и маршрутизатора. Storage-узел отвечает за физическую таблицу или shard таблицы. Registry управляет регистрацией узлов хранения и запуском процесса \texttt{qdb-storage}; Monitor выполняет периодические проверки доступности через внутреннее действие \texttt{ping}.

При недоступности Storage-узла несколько последовательных heartbeat-проверок завершаются ошибкой, после чего узел помечается как недоступный и может быть перезапущен. Это решение не заменяет полноценную репликацию, но закрывает требование автоматического обнаружения и перезапуска отказавшего узла.

Для управления таблицами Entrypoint использует каталог. Каталог хранит логическую структуру баз данных и таблиц, а Registry связывает запись каталога с конкретным процессом Storage. При создании таблицы сервер сохраняет схему, выделяет каталог для shard-файлов и поднимает соответствующий Storage-процесс. При удалении таблицы запись каталога удаляется, а связанные файлы и индексные структуры очищаются.

\section{Язык запросов и обработка команд}

Лексер разбивает входной текст на токены: ключевые слова, идентификаторы, строковые литералы, числовые литералы, операторы и знаки пунктуации. Ключевые слова регистронезависимы, идентификаторы сохраняют исходный регистр. Строковые литералы заключаются в двойные кавычки.

Парсер строит AST для команд управления базами данных, DDL, DML, команд доступа и темпорального отката. Выражения условий представлены отдельной иерархией узлов: сравнение, \texttt{BETWEEN}, \texttt{LIKE}, \texttt{AND}, \texttt{OR}. Для агрегатов используются узлы \texttt{AggregateExpr} с типами \texttt{SUM}, \texttt{COUNT}, \texttt{AVG}.

Парсер реализован как рекурсивный спуск. Для каждого уровня грамматики выделена отдельная функция: разбор оператора верхнего уровня, табличной ссылки, литерала, выражения, списка колонок, условия и timestamp для \texttt{REVERT}. Приоритет логических операций задается структурой грамматики: сначала разбираются предикаты, затем цепочки \texttt{AND}, затем цепочки \texttt{OR}. Это позволяет корректно обрабатывать выражения вида \texttt{a == 1 OR b == 2 AND c == 3} без дополнительной перестановки AST.

\begin{lstlisting}[style=qdbcpp,caption={Фрагмент иерархии условий AST}]
struct Condition : ASTNode {
    enum class Kind { Comparison, Between, Like, And, Or };
    virtual auto KindOf() const -> Kind = 0;
};

struct ComparisonCondition : Condition {
    std::unique_ptr<Expression> Left;
    Operator Op;
    std::unique_ptr<Expression> Right;
};
\end{lstlisting}

Семантический анализ проверяет существование таблиц и колонок, соответствие типов, допустимость \texttt{NULL} для \texttt{NOT\_NULL}-полей, наличие значений по умолчанию и права пользователя. После этого сервер формирует внутреннее представление запроса и передает его на Storage.

Ошибки разбора и выполнения не приводят к аварийному завершению программы. Компоненты выбрасывают исключения с диагностическим текстом, а верхний уровень обработки преобразует их в JSON-ответ со статусом \texttt{error}. Благодаря этому клиент получает информативное описание проблемы: неизвестная колонка, нарушение количества значений, неверный тип литерала, отсутствие прав или ошибка доступа к файлу.

\section{Физическое хранение данных}

Подсистема хранения размещается в \texttt{apps/storage}. Основной объект хранения — таблица. Для каждой таблицы создается набор файлов: файл страниц данных, файл схемы, файл интернера строк, журнал изменений и индексные файлы.

\begin{table}[H]
\caption{Файлы одной таблицы}
\centering
\small
\begin{tabularx}{\textwidth}{>{\raggedright\arraybackslash}p{0.36\textwidth}Y}
\toprule
Файл & Назначение \\
\midrule
\breaktt{<table>.data} & Страничное хранение записей и метаданных таблицы. \\
\breaktt{<table>.schema} & Сериализованная схема таблицы: колонки, типы, флаги, значения по умолчанию. \\
\breaktt{<table>.bin} & Файл интернера со строковыми значениями (\texttt{StringId} и строка). \\
\breaktt{<table>.jnl} & Журнал обратных операций для темпорального отката. \\
\breaktt{<table>\_id\_to\_addr.idx} & Индекс соответствия идентификатора записи и физического адреса. \\
\breaktt{<table>\_<column>.idx} & Индекс для пользовательской колонки с модификатором \texttt{INDEXED}. \\
\bottomrule
\end{tabularx}
\end{table}

\subsection{Страничный менеджер}

Pager абстрагирует файловый ввод-вывод и работает страницами фиксированного размера 4096 байт. Он поддерживает чтение страницы по номеру, запись страницы, добавление новой страницы и усечение файла. Такой слой упрощает реализацию таблиц и индексов: все структуры хранятся как последовательность страниц одинакового размера.

Внутри страницы таблицы используется последовательная структура записей. Заголовок страницы хранит идентификатор страницы, количество записей и смещение следующего свободного байта. Каждая запись представлена ячейкой с идентификатором, размером и флагом удаления, за которой следует сериализованный полезный блок. Новые записи добавляются в конец страницы; освобождаемое место от удаленных записей не компактизируется, поэтому учитывается только непрерывный свободный хвост.

Алгоритм вставки записи сначала сериализует значения согласно схеме, затем ищет страницу, где свободного хвоста достаточно для ячейки и данных. Если такая страница найдена, запись добавляется в конец; иначе создается новая страница. При чтении запись извлекается по номеру слота, а удаленные записи пропускаются.

\subsection{Записи и значения}

Значение представлено вариантом из трех типов: \texttt{NULL}, \texttt{int32\_t} и интернированная строка (\texttt{StringId}). При сериализации каждому значению предшествует байт типа, далее хранится либо целое число, либо идентификатор строки. Строковые данные физически размещаются в файле \texttt{.bin} интернера, а записи содержат только ссылку на них. Запись адресуется парой \texttt{page\_idx} и \texttt{slot\_idx}.

Схема таблицы хранит имена колонок, типы \texttt{INT} и \texttt{STRING}, флаги \texttt{NOT\_NULL} и \texttt{INDEXED}, а также значение \texttt{DEFAULT}. Для индексируемой колонки автоматически включается \texttt{NOT\_NULL}, что синхронизирует уникальные индексы и требования целостности.

При вставке с явным списком колонок таблица сначала создает запись размера схемы, затем применяет значения по умолчанию для непереданных столбцов и только после этого записывает переданные пользователем значения. Если обязательная колонка остается пустой, операция завершается семантической ошибкой. При обновлении старая версия записи читается по физическому адресу, журналируется, удаляется из страницы и заменяется новой версией; индексы синхронно обновляют ссылки на актуальный адрес.

\section{Индексирование на основе B*+-дерева}

Вариант проекта использует B*+-дерево. Индекс хранит только ключи и физические адреса записей, что соответствует требованию о запрете дублирования самих данных. Для каждой индексируемой колонки создается отдельный индексный файл, а служебный индекс \texttt{id\_to\_addr} ускоряет чтение записи по внутреннему идентификатору.

B*+-дерево отличается от классического B+-дерева тем, что при переполнении сначала пытается перераспределить ключи между соседними узлами. Если перераспределение невозможно, два узла разделяются на три, что поддерживает большую минимальную заполненность страниц. Узел дерева занимает одну страницу размером 4096 байт. Листовые узлы связаны списком, поэтому диапазонные запросы выполняются последовательным проходом по листьям после поиска первого подходящего ключа.

\begin{figure}[H]
\centering
\fbox{\begin{minipage}{0.88\textwidth}
\centering
Root/Internal node\\
$\downarrow$ поиск дочернего диапазона\\
Leaf node $\rightarrow$ Leaf node $\rightarrow$ Leaf node\\
ключи листьев содержат адреса записей в \texttt{.data}
\end{minipage}}
\caption{Логическая структура B*+-индекса}
\end{figure}

Для строковых индексов используется ключ \texttt{StringId}, получаемый через интернер. Это уменьшает размер ключей и упрощает сравнение в дереве.

Исполнитель запросов использует индекс, когда условие содержит равенство по индексируемой колонке и сравнение можно свести к паре «колонка — литерал». Для составного условия \texttt{AND} исполнитель пытается найти индексируемый предикат слева или справа, получает кандидатов по индексу, а затем повторно применяет полное условие к найденным записям. Если подходящего индекса нет, выполняется последовательный просмотр записей таблицы. Такой подход сохраняет корректность для сложных условий и ускоряет типичный случай точечного поиска.

\begin{table}[H]
\caption{Сложность основных операций индекса}
\centering
\small
\begin{tabularx}{\textwidth}{>{\raggedright\arraybackslash}p{0.28\textwidth}Y}
\toprule
Операция & Оценка и пояснение \\
\midrule
Поиск ключа & \texttt{O(log N)} по высоте дерева, затем последовательный проход по совпадающим ключам в листовых узлах. \\
Вставка & \texttt{O(log N)} на поиск листа; при переполнении выполняется перераспределение или split с обновлением родителя. \\
Удаление & \texttt{O(log N)} на поиск листа; при недозаполнении выполняется заимствование у соседей или слияние. \\
Диапазон & \texttt{O(log N + K)}, где \texttt{K} — число записей в диапазоне, потому что листья связаны последовательным списком. \\
\bottomrule
\end{tabularx}
\end{table}

\section{Журналирование и темпоральный откат}

Для реализации команды \texttt{REVERT} используется журнал обратных операций. Перед изменением таблицы в файл \texttt{.jnl} записывается операция, позволяющая отменить предстоящее изменение. Например, перед вставкой записи журнал содержит обратную операцию удаления, перед удалением — старую версию записи для повторной вставки, перед обновлением — предыдущую версию записи.

Каждая запись журнала содержит timestamp в формате \texttt{yyyy.mm.dd-hh:mm:ss.mmm}, идентификатор записи, тип обратной операции, сериализованные данные при необходимости и размер записи журнала. Последний размер позволяет читать журнал с конца к началу без полного сканирования файла.

\begin{lstlisting}[style=qdbcpp,caption={Логика чтения журнала при откате}]
while (position > 0) {
    read_last_track_size();
    seek_to_track_begin();
    read_time_record_id_and_type();

    if (track_time < target_time) {
        break;
    }

    apply_inverse_operation();
    position = current_track_begin;
}
\end{lstlisting}

Полные снимки файлов не создаются. Благодаря этому расход дискового пространства зависит от числа изменений, а не от полного объема базы. После отката новая запись в таблицу усекает ветку журнала после текущей позиции, что предотвращает применение уже неактуальной истории.

Порядок записи выбран так, чтобы журнал появлялся на диске раньше изменения таблицы. Для операции вставки сначала сохраняется обратная операция удаления, затем увеличивается счетчик идентификаторов, обновляется схема, записывается страница данных и обновляются индексы. Для удаления, наоборот, журнал содержит полную старую запись, чтобы откат мог вернуть ее в таблицу и восстановить индексные ссылки. Такое журналирование не является полноценной транзакционной моделью с конкурирующими транзакциями, но выполняет требование темпорального отката без snapshot-копирования файлов.

\section{Интернирование строк}

Для оптимизации оперативной памяти используется Interner. Он хранит каждую уникальную строку в единственном экземпляре, а записи ссылаются на уже существующее представление. Внутри применяется \texttt{std::deque<std::string>} для стабильности адресов строк и две таблицы сопоставления: \texttt{std::unordered\_map<std::string, StringId>} и \texttt{std::unordered\_map<uint32\_t, std::string\_view>}.

Интернирование особенно полезно для повторяющихся значений в таблицах: названий категорий, статусов, городов, групп пользователей. Кроме экономии памяти, оно ускоряет некоторые операции сравнения, поскольку повторяющиеся строки имеют общие представления.

Пул строк связан с временем жизни базы данных и использует файл \texttt{.bin} для хранения строк в виде пар (идентификатор, строка). При чтении записи значение восстанавливается по \texttt{StringId}, поэтому одинаковые значения, пришедшие из разных страниц, получают единое представление в памяти и доступ через \texttt{string\_view}.

\section{Клиентское приложение}

Клиент поддерживает интерактивный и пакетный режимы. В интерактивном режиме пользователь запускает программу без аргументов и вводит команды в терминале. В пакетном режиме клиент получает путь к файлу сценария и последовательно отправляет команды из него. В обоих режимах команда считается завершенной после символа \texttt{;}.

Ответ сервера разбирается как JSON. Успешный \texttt{SELECT} отображается как массив объектов. Ошибки выводятся с человекочитаемым сообщением, полученным из поля \texttt{message}. Если сервер возвращает статус \texttt{pending}, клиент может обращаться к API проверки задачи и ожидать завершения длительной операции.

Клиентская часть отделяет чтение команд от их отображения. Reader отвечает за интерактивный ввод и чтение сценариев из файла, а Renderer отвечает за форматирование результата. Такое разделение упрощает тестирование: можно отдельно проверять многострочный ввод, пакетный режим, обработку пустых команд, вывод успешного результата и вывод ошибок.

\section{Безопасность и разграничение доступа}

Подсистема безопасности включает хранилище аккаунтов, хеширование паролей с солью, выпуск JWT-токенов и проверку прав доступа. Пароли не сохраняются в открытом виде: для вычисления хеша используется PBKDF2-HMAC-SHA256 с большим числом итераций. Сравнение хеша выполняется через константное по времени сравнение, что снижает риск утечки информации по времени ответа.

JWT содержит subject пользователя, время выпуска и срок действия. После успешной аутентификации клиент передает токен в последующих запросах. Сервер проверяет подпись и срок действия токена, затем применяет RBAC-правила к операции.

Модель доступа учитывает личные права пользователя и правила по умолчанию через wildcard-записи. Для операций различаются чтение, запись, создание и удаление.

Правила доступа сериализуются в JSON-файл как набор записей с пользователем, базой, таблицей и множеством разрешений. Для универсальных правил используется символ \texttt{*}: он позволяет описать права по умолчанию на все базы или все таблицы. Проверка права выполняется поиском правила, у которого совпадает пользователь либо wildcard, совпадает база либо wildcard, совпадает таблица либо wildcard и содержится требуемое разрешение.

\section{Логирование, телеметрия и асинхронные задачи}

Access log фиксирует тело запроса, идентификатор клиента, идентификатор обработчика, время начала и завершения обработки, а также статус результата. Запись логов вынесена из основной логики обработки запроса, чтобы снизить влияние дискового ввода-вывода на задержку ответа.

Телеметрия собирает текущий RPS, средний и максимальный RPS за десятиминутное окно, среднюю длительность обработки за десять секунд и долю ошибок за последнюю минуту. Для счетчиков и оконных структур используются потокобезопасные механизмы, в том числе атомарные счетчики там, где это возможно без потери корректности.

Длительные запросы могут выполняться асинхронно. Entrypoint создает GUID v4, помещает задачу в потокобезопасную очередь и сразу возвращает клиенту статус \texttt{pending}. Worker-потоки извлекают задачи, отправляют их на Storage и сохраняют итоговый результат. Клиент получает результат через запрос \texttt{check\_task}.

Для телеметрии используются скользящие окна разной длины: текущий RPS считается по последней завершенной секунде, средний и максимальный RPS — по десятиминутному окну, средняя длительность — по десятисекундному окну, а доля ошибок — по минутному окну. Это позволяет одновременно видеть мгновенную нагрузку и сглаженную картину поведения сервера.

\section{Тестирование}

Проект использует GTest. Тестами покрываются лексер, парсер, сетевые примитивы, клиентские компоненты, pager, schema, journal, value, interner, таблицы, B*+-дерево, executor, RBAC, аутентификация, телеметрия и асинхронный статус. Сборка организована через CMake и Conan, а вспомогательные команды вынесены в Makefile.

Для контроля качества применяются clang-format и clang-tidy. Разделение на модули и наличие unit-тестов позволяют проверять изменения локально и поддерживать стабильность интерфейсов между клиентом, сервером и хранилищем.

\begin{table}[H]
\caption{Основные группы тестов}
\centering
\small
\begin{tabularx}{\textwidth}{>{\raggedright\arraybackslash}p{0.24\textwidth}Y}
\toprule
Подсистема & Проверяемые свойства \\
\midrule
Parser/Lexer & Токенизация, построение AST, поддержка SQL-конструкций и диагностика ошибочных запросов. \\
Storage & Страницы, схема, записи, индексы, журнал, агрегаты, SQL-логика и восстановление состояния. \\
Server & API, маршрутизация, безопасность, RBAC, телеметрия, асинхронные статусы и обработка ошибок. \\
Client & Чтение команд, конфигурация, интерактивный и пакетный режимы, форматирование ответов. \\
Core & TCP-клиент, TCP-сервер, сокеты, JSON-файлы, thread pool и интеграционные сетевые сценарии. \\
\bottomrule
\end{tabularx}
\end{table}

Проверка сборки выполняется командами \texttt{make build} и \texttt{make test}. Для локального анализа качества кода предусмотрена команда \texttt{make check}, объединяющая проверки форматирования и статического анализа. Такой набор команд делает результат воспроизводимым для всех участников команды.

\section*{Вывод}
\addcontentsline{toc}{section}{Вывод}

В ходе курсовой работы разработана СУБД QuasarDB на C++20. Проект включает клиентский терминал, входной сервер, узлы хранения и общую библиотеку сетевых примитивов. Реализованные решения соответствуют основным требованиям задания: данные сохраняются в файловой системе, запросы проходят разбор и валидацию, результаты возвращаются в JSON, индексируемые поля используют B*+-дерево без дублирования записей, а изменения журналируются для последующего темпорального отката.

Ключевыми техническими результатами стали страничный движок хранения, последовательная организация записей, сериализация схем, интернирование строк, индексные файлы, журнал обратных операций, SQL-подобный язык с составными условиями и агрегатами, клиент-серверный обмен JSON-сообщениями, аутентификация через JWT, RBAC, логирование и телеметрия. Эти подсистемы демонстрируют практическое применение системного программирования, структур данных, многопоточности и файлового ввода-вывода.

\clearpage
\section*{Список использованных источников}
\addcontentsline{toc}{section}{Список использованных источников}

\begin{enumerate}[label={\arabic*.}]
    \item ГОСТ 7.32--2017. Система стандартов по информации, библиотечному и издательскому делу. Отчет о научно-исследовательской работе. Структура и правила оформления. Москва: Стандартинформ, 2017.
    \item ГОСТ Р 7.0.100--2018. Система стандартов по информации, библиотечному и издательскому делу. Библиографическая запись. Библиографическое описание. Общие требования и правила составления. Москва: Стандартинформ, 2018.
    \item ISO/IEC 14882:2020. Programming languages -- C++. Geneva: International Organization for Standardization, 2020.
    \item CMake Documentation. URL: \url{https://cmake.org/documentation/} (дата обращения: 27.05.2026).
    \item nlohmann/json: JSON for Modern C++. URL: \url{https://github.com/nlohmann/json} (дата обращения: 27.05.2026).
    \item GoogleTest User's Guide. URL: \url{https://google.github.io/googletest/} (дата обращения: 27.05.2026).
    \item Comer D. The Ubiquitous B-Tree. ACM Computing Surveys. 1979. Vol. 11, No. 2. P. 121--137.
    \item Bayer R., McCreight E. Organization and Maintenance of Large Ordered Indexes. Acta Informatica. 1972. Vol. 1. P. 173--189.
\end{enumerate}

\clearpage
\appendix
\section{Грамматика основных команд}

\begin{lstlisting}[style=qdbcpp,caption={Фрагмент EBNF-грамматики языка QuasarDB}]
<statement> ::= ( <db_stmt> | <ddl_stmt> | <dml_stmt> ) ";"

<create_db> ::= "CREATE" "DATABASE" <identifier>
<drop_db>   ::= "DROP"   "DATABASE" <identifier>
<use_db>    ::= "USE"               <identifier>

<create_table> ::= "CREATE" "TABLE" <table_ref>
                   "(" <col_def> { "," <col_def> } ")"

<insert_stmt> ::= "INSERT" "INTO" <table_ref>
                  "(" <identifier> { "," <identifier> } ")"
                  "VALUE" <row_values> { "," <row_values> }

<select_stmt> ::= "SELECT" <select_list>
                  "FROM" <table_ref>
                  [ "WHERE" <condition> ]
\end{lstlisting}

\section{Пример сценария работы}

\begin{lstlisting}[style=qdbcpp,caption={Пример SQL-сценария для QuasarDB}]
CREATE DATABASE shop;
USE shop;

CREATE TABLE products (
    id INT NOT_NULL INDEXED DEFAULT 0,
    name STRING NOT_NULL,
    price INT DEFAULT 0
);

INSERT INTO products (id, name, price)
VALUE (1, "Apple", 99),
      (2, "Banana", 75);

SELECT name AS product_name, AVG(price) AS avg_price
FROM products
WHERE price >= 1 AND name LIKE "^A.*";
\end{lstlisting}

\section{Структура репозитория}

\begin{lstlisting}[style=qdbcpp,caption={Основные каталоги проекта}]
apps/client   - CLI client terminal
apps/server   - Entrypoint server: routing, security, observability
apps/storage  - Storage node: data, indexes, journal
libs/core     - Shared code and network primitives
docs          - Documentation, grammar, API and report
external      - Conan dependency manifests
\end{lstlisting}

\end{document}
